\chapter{Geografía y Distritos de NeoBilbao}

\neofigure{01_distrito7_la_margen.jpg}{Vista urbana de La Margen: asfalto empapado, neón parpadeante y viaductos del tren mag-lev.}

La metrópolis de NeoBilbao está encajada verticalmente en una gigantesca falla tectónica del planeta Ur-Bare. Desde las profundidades incandescentes de la Ría de Acero hasta la cúspide bioclimática de la Cúpula Dorada, la ciudad es una manifestación física de la jerarquía social corporativa.

\section{1. La Margen (Distrito 7, Nivel Medio)}

Dominante durante la era de 2084 y todavía existente como sector residencial denso en 2102, \textbf{La Margen} es el hogar de la clase trabajadora urbana, técnicos de grado medio, auxiliares y empleados de mantenimiento.

\begin{itemize}
  \item \lead{Los bloques colmena:} Rascacielos modulares de hormigón poroso y vigas de acero oxidado. Los ciudadanos habitan en micro-apartamentos cápsula de doce a catorce metros cuadrados con literas abatibles, respiradores murales de cuota y cocinas sintéticas compartidas por planta.
  \item \lead{La red Mag-Lev (Línea 4):} Monorrieles elevados de tracción magnética cuyos viaductos cruzan sobre las calles a la altura del tercer piso. El traqueteo sordo y el chirrido de los electroimanes sacuden los edificios cada diez minutos.
  \item \lead{La cultura de las tascas:} Lugares de encuentro popular situados bajo los puentes y en los callejones resguardados de la lluvia. En locales como el bar «Zulo», los obreros beben kalimotxo sintético o cerveza de malta artificial, intercambian herramientas clandestinas y se desahogan lejos de los micrófonos corporativos.
  \item \lead{Los Dispensarios Periféricos:} Centros sanitarios masificados como el de Gurutzeta Periférico, donde médicos y celadores atienden a destajo quemaduras químicas, asma por sirimiri y traumatismos industriales con presupuestos bajo mínimos.
\end{itemize}

\section{2. Los Zócalos Industriales (El Abismo Inferior)}

\neofigure{foundry_workers_platform.jpg}{Operarios sobre las pasarelas metálicas resbaladizas en los Zócalos de fundición.}

El corazón productivo de la colonia en 2102. Es una fosa colosal donde el sol nunca penetra y donde la temperatura roza los cincuenta grados centígrados.

\begin{itemize}
  \item \lead{La Ría de Acero:} El lecho de la garganta tectónica por donde fluyen ríos incandescentes de escoria y metal fundido. El resplandor rojizo y el siseo del vapor químico crean una atmósfera asfixiante y perpetuamente encendida.
  \item \lead{Los Altos Hornos de Ur-Bare:} Titánicas estructuras siderúrgicas levantadas contra las paredes del cañón. Aquí se refina el Burdin-Urdin mediante reactores térmicos y hornos de plasma que operan día y noche en ciclos ininterrumpidos.
  \item \lead{El Puente Transbordador:} La icónica barquilla de carga suspendida por cables de alta tensión sobre el abismo. Transporta contenedores pesados de mineral y cuadrillas enteras de mineros entre los dos extremos del desfiladero.
  \item \lead{La Aguja:} Una colosal torre vertical de ascensores hidráulicos y conductos de ventilación que conecta directamente los niveles más bajos de la fundición con las plantas de procesado y los accesos restringidos de la zona superior.
  \item \lead{Canales de refrigeración y colectores:} Un laberinto sumergido de tuberías de desagüe de tres a seis metros de diámetro donde se evacúa el agua radiactiva y los efluentes ácidos hacia las fosas exteriores.
\end{itemize}

\section{3. La Cúpula Dorada (La Cúspide Corporativa)}

\neofigure{elites_golden_dome.jpg}{La élite corporativa paseando bajo la luz solar simulada y los jardines de la Cúpula Dorada.}

El paraíso cerrado reservado para los accionistas, altos ejecutivos y familias dinásticas de Orekak-Corp (\emph{Zaharrak}).

\begin{itemize}
  \item \lead{El escudo bioclimático:} Una cúpula geodésica de polímero transparente y aleaciones de titanio que filtra la radiación cósmica y proyecta ciclos diurnos de sol dorado y cielos azules simulados.
  \item \lead{Arquitectura y jardines:} Torres sinuosas de líneas orgánicas inspiradas en el titanio del museo Guggenheim, rodeadas de parques con hierba real, fuentes de agua dulce filtrada y aire con aromas cítricos y pino.
  \item \lead{El perímetro de exclusión:} El acceso desde los niveles inferiores está protegido por esclusas de descompresión blindadas, campos láser biométricos y puestos de guardia militarizados con orden de disparar a discreción ante cualquier intrusión no autorizada.
\end{itemize}

\section{4. El Yermo Exterior de Ur-Bare}

Más allá del perímetro de contención de la colonia se extiende la desolación de Ur-Bare:
\begin{itemize}
  \item \lead{Tormentas de arena ferrosa:} Vientos huracanados que arrastran polvo de óxido cortante como cuchillas, capaz de erosionar trajes de presión en minutos.
  \item \lead{Las Minas Abiertas de Superficie:} Canteras abandonadas donde los primeros colonos excavaron a cielo abierto antes de refugiarse en la Ría.
\end{itemize}
